\begin{definition}[Funções $C^{1}$]
  Uma função $f:U\subseteq\mathbb{R}^{n}\to\mathbb{R}$ é de classe $C^{1}$ e escrevemos $f\in C^{1}$ se 
  \begin{align*}
    \frac{\partial f}{\partial x_i}:U&\to\mathbb{R},\\
    x&\to\frac{\partial f}{\partial x_i}(x),
  \end{align*}
  é continua em $U$ para tudo $i=1,\cdots,n$. 
\end{definition}
\begin{theorem}
  Tudo função $f\in C^{1}$ é diferenciável. 
\end{theorem}
\begin{proof} 
  Sem perda de generalidade seja $n=2$.

  Fixemos $c=(a,b)\in\mathbb{R}^{2}$ e $v=(h,k)\in\mathbb{R}^{2}$ tal que $c+v\in U$.

  Se definirmos $r$ e usarmos o Teorema do Valor Médio em uma dimenção, temos que existem $\theta_1,\theta_2\in[0,1]$ tais que 
  \begin{align*}
    r(v)=r(h,k)&=f(a+h,b+k)-f(a,b)-\frac{\partial f}{\partial x}(a,b)h-\frac{\partial f}{\partial y}(a,b)k\\
    &=f(a+h,b+k)-f(a,b+k)+f(a,b+k)-f(a,b)-\frac{\partial f}{\partial x}(a,b)h-\frac{\partial f}{\partial y}(a,b)k\\
    &=\frac{\partial f}{\partial x}(a+\theta_1h,b+k)h+\frac{\partial f}{\partial y}(a,b+\theta_2k)k-\frac{\partial f}{\partial x}(a,b)h-\frac{\partial f}{\partial y}(a,b)k\\
    &=\left[ \frac{\partial f}{\partial x}(a+\theta_1h,b+k)-\frac{\partial f}{\partial x}(a,b) \right]h+\left[ \frac{\partial f}{\partial y}(a,b+\theta_2k)-\frac{\partial f}{\partial y}(a,b) \right]k.
  \end{align*}
  Assim,
  \begin{align*}
    \left|\frac{r(v)}{|v|}\right|&\leq\frac{|h|}{\sqrt{h^2+k^2}}\left| \frac{\partial f}{\partial x}(a+\theta_1h,b+k)-\frac{\partial f}{\partial x}(a,b) \right|+\frac{|k|}{\sqrt{h^2+k^2}}\left| \frac{\partial f}{\partial y}(a,b+\theta_2k)-\frac{\partial f}{\partial y}(a,b) \right|\\
    &\leq\left| \frac{\partial f}{\partial x}(a+\theta_1h,b+k)-\frac{\partial f}{\partial x}(a,b) \right|+\left| \frac{\partial f}{\partial y}(a,b+\theta_2k)-\frac{\partial f}{\partial y}(a,b) \right|.
  \end{align*}
  Logo
  \begin{align*}
    \lim_{v \to 0}\left| \frac{r(v)}{|v|} \right|\leq \lim_{v \to 0}\left| \frac{\partial f}{\partial x}(a+\theta_1h,b+k)-\frac{\partial f}{\partial x}(a,b) \right|+\left| \frac{\partial f}{\partial y}(a,b+\theta_2k)-\frac{\partial f}{\partial y}(a,b) \right|=0.
  \end{align*}
  O que nos permite concluir que
  \begin{align*}
    \lim_{v \to 0}\frac{r(v)}{|v|}=0,
  \end{align*}
  i,e, $f$ é uma função diferenciável. 
\end{proof}
\newpage
\begin{theorem}
  Seja $U\subseteq \mathbb{R}^{m}$, $V\subseteq\mathbb{R}^{n}$ e $f:U\to V$ uma aplicação cujas funções coordenadas $f_1,\cdots,f_n:U\to \mathbb{R}$ possuem derivadas parciais em $a\in U$ e $g:V\to\mathbb{R}$ uma função diferenciável no ponto $b=f(a)$.\\
  Então $g\circ f$ possui derivada parciais no ponto $a$ e vale
  \begin{align*}
    \frac{\partial (g\circ f)}{\partial x_i}(a)=\sum_{k=1}^{n}\frac{\partial g}{\partial y_k}(b)\cdot\frac{\partial f_k}{\partial x_i}(a),\quad i=1,\cdots,m.
  \end{align*}
  Além disso, $f,g\in C^1$ implicam que $g\circ f\in C^1$. 
\end{theorem}
\begin{proof}
  Tarefa...
\end{proof}
\begin{corollary}
  Seja $f:U\to \mathbb{R}$ diferenciável no ponto aberto $U\subseteq \mathbb{R}^{n}$, com $a\in U$. Dadoo o vetor $v=(\alpha_1,\cdots,\alpha_n)\in\mathbb{R}^{n}$ se $\lambda:(-\delta,\delta)\to U$ é qualquer caminho diferenciável tal que $\lambda(0)=a$ é $\lambda^{\prime}(0)=v$, tem-se
  \begin{align*}
    (f\circ \lambda)^{\prime}(0)&=\langle \grad f(a),v\rangle\\
    &=\frac{\partial f}{\partial v}(a)\\
    &=\sum_{i=1}^{n}\frac{\partial f}{\partial x_i}(a)\alpha_i.
  \end{align*}
\end{corollary}
\begin{corollary}[T.V.M]
  Dada uma $f:U\to\mathbb{R}$ diferenciável no aberto $U\subseteq \mathbb{R}^{n}$, se o segmento de reta $[a,a+v]\subseteq U$ então existe $\theta\in(0,1)$ tal que
  \begin{align*}
    f(a+v)-f(a)&=\frac{\partial f}{\partial v}(a+\theta v)\\
    &=\langle \grad f(a+\theta v),v \rangle\\
    &=\sum_{i=1}^{n}\frac{\partial f}{\partial x_i}(a+\theta v)\alpha_i.
  \end{align*} 
\end{corollary}
\begin{proof} 
  Considere o caminho $\lambda:[0,1]\to\mathbb{R}$ dado por $\lambda(t):a+tv$, é evidente que $\lambda(0)=a$ e $\lambda(1)=v$, logo usando o teorema de valor medio no una dimenção ten-se que
  \begin{align*}
    f(a+v)-f(a)&=f(\lambda(1))-f(\lambda(0))\\
    &=(f\circ \lambda)^{\prime}(\theta)\\
    &=\sum_{i=1}^{n}\frac{\partial f}{\partial x_i}(a+\theta v)\alpha_i\\
    &=\frac{\partial f}{\partial v}(a+\theta v).
  \end{align*}
\end{proof}
\begin{corollary}
  Seja $f:U\to\mathbb{R}$ diferenciável no aberto $U\subseteq \mathbb{R}^{n}$. Se $U$ é conexo e $\grad f(x)=0$ para tudo $x\in U$, então é constante. 
\end{corollary}
\begin{proof} 
  Pela regra da cadeira resulta que
  \begin{align*}
    \frac{\partial f}{\partial v}(a)=\langle \grad f(a),v\rangle=0.
  \end{align*}
\end{proof}
\begin{definition}[Classe $C^{k}$]
  Dizer que $f^{\prime}:U\to L(\mathbb{R}^{m},\mathbb{R}^{n})$ é contínua equivale a afirmar que cada parcial $\frac{\partial f_i}{\partial x_j}:U\to\mathbb{R}$ é continua, i,e, que $f\in C^{1}$.

  Por indução define-se classe $C^{k}$ por $f\in C^{k}$ si é diferenciável e sua derivada $f^{\prime}:U\to\mathbb{R}$ é de classe $C^{k-1}$.

  Si $f\in C^{k}$ para tudo $k$, diz-se que $f\in C^{\infty}$
\end{definition}
\begin{proposition}
  Seja $f:U\subseteq\mathbb{R}^{m}\to\mathbb{R}$ uma função que possui todas as derivadas parciais $\frac{\partial f}{\partial x_i}(a)$ com $i=1,\cdots,n$. Então, em cada direção $e_i$, a restrição de $f$ no segmento da reta $L_i=\{a+te_i:|t|<\delta\}$ é contínua no ponto $a$.

  No entanto isso não garante que $f$ seja contínua en $a$ como função de varias variaveis. 
\end{proposition}
\begin{proof} 
  Para um $i$ fixo defina
  \begin{align*}
    \phi_i(t)=f(a+te_i),\quad t\in(-\delta,\delta),
  \end{align*}
  onde $\delta>0$ é suficientemente pequeno.

  Como a derivada $\frac{\partial f}{\partial x_i}(a)$ existe, então
  \begin{align*}
    \frac{\partial f}{\partial x_i}(a)=\lim_{t \to 0}\frac{\phi_i(t)-\phi_i(0)}{t}=\phi_i^{\prime}(0).
  \end{align*}
  Então, $\phi_i$ é derivavel em $t=0$. Logo $\phi_i$ é continua em $t=0$, i,e
  \begin{align*}
    \lim_{t \to 0}f(a+te_i)=f(a).
  \end{align*}
\end{proof}
\begin{example}
  Considere
  \begin{align*}
    f:\mathbb{R}^{2}&\to\mathbb{R},\\
    f(x,y)&\to\begin{cases}
      1, &\text{ se } xy\neq 0 \text{,} \\
      0, &\text{ se } xy=0 .
    \end{cases}
  \end{align*}
  Observe que
  \begin{align*}
    \frac{\partial f}{\partial x}(0,0)=0,\quad\frac{\partial f}{\partial y}(0,0)=0.
  \end{align*}
  Pelo
  \begin{align*}
    \lim_{(x,y) \to (0,0)}f(x,y)=\lim_{t \to 0}f(t,t)=&1,\\
    \lim_{(x,y) \to (0,0)}f(x,y)=\lim_{t \to 0}f(0,t)=&0.
  \end{align*}
  Então $f$ no é contínua. 
\end{example}
\begin{proposition}
  Seja $f:U\subseteq\mathbb{R}^{m}\to\mathbb{R}$ uma função que possui derivada parcial $\frac{\partial f}{\partial x_i}(x)$ para tudo $x$ no segmento da reta $J_i=[a-\delta e_i,a+\delta e_i]$ com $\delta>0$. Se $\frac{\partial f}{\partial x_i}(x)>0$ para tudo $x\in J_i$, então $f$ é estritamente crescente ao longo do segmento $J_i$ na direção $e_i$, ou seja se $s<t$, então $f(a+se_i)<f(a+te_i)$ para tuddo $s$ e $t$ com $|s|,|t|<\delta$. 
\end{proposition}
\begin{proof} 
  Defina
  \begin{align*}
    \phi_i(t)=f(a+te_i),\quad |t|<\delta.
  \end{align*}
  Logo
  \begin{align*}
    \phi^{\prime}(t)=\frac{\partial f}{\partial x_i}(a+te_i)>0.
  \end{align*}
  Então, si voce toma $t-s>0$ (i,e, $t>s$) ten-se que por ele teorema de valor médio
  \begin{align*}
    \phi(t)-\phi(s)=\phi^{\prime}(\theta)(t-s)>0.
  \end{align*}
  Isso es
  \begin{align*}
    f(a+te_i)>f(a+se_i).
  \end{align*}
\end{proof}
\begin{example}[Transformações lineares]
  Se $T:\mathbb{R}^{m}\to\mathbb{R}^{n}$ é  uma transformação linear, então $T$ é diferenciável, pois
  \begin{align*}
    T(x+v)-T(x)=T(v)+0.
  \end{align*}
  Entáo $T^{\prime}(x)=T$ (constante) , logo $T^{(n)}=0$ com $n\geq 2$. 
\end{example}
\begin{example}[Soma dos componentes]
  Seja
  \begin{align*}
    S:\mathbb{R}^{m}\times\mathbb{R}^{m}&\to\mathbb{R}^{m}\\
    (x,y)&\to x+y.
  \end{align*}
  Então
  \begin{align*}
    S(x+u,y+v)=S(x,y)+S(x,v)+S(u,y)+S(u,v).
  \end{align*}
  Logo
  \begin{align*}
    S(x+u,y+v)-S(x,y)=S(x,v)+S(u,y)+S(u,v).
  \end{align*}
  Então
  \begin{align*}
    S^{\prime}(x,y)\cdot(u,v)=S(u,v)=u+v.
  \end{align*}
\end{example}
\begin{example}[Aplicações bilineares]\label{ex:aplicacoes-bilineares}
  Seja $B:\mathbb{R}^{m}\times\mathbb{R}^{n}\to\mathbb{R}^{p}$ uma aplicação bilinear, ou seja
  \begin{itemize}
    \item $y$ fixo, $x\to B(x,y)$ é linear. 
    \item $x$ fixo, $y\to B(x,y)$ é linear. 
  \end{itemize}
  Tomemos as bases canonicas
  \begin{itemize}
    \item $\{e_1,\cdots,e_m\}\in\mathbb{R}^{m}$.
    \item $\{f_1,\cdots,f_n\}\in\mathbb{R}^{n}$.
  \end{itemize}
  Definimos os vetores $v_{ij}=B(e_i,f_j)\in\mathbb{R}^{p}$, então usando a bilinearidade
  \begin{align*}
    B(x,y)&=B\left( \sum_{i=1}^{m}x_ie_i,\sum_{j=1}^{m}y_jf_j \right),\\
    &=\sum_{i=1}^{m}\sum_{j=1}^{n}x_iy_jv_{ij}.
  \end{align*}
  Logo $B$ é um polinomio has coordenadas $x,y$, assim $B$ é contínua.\\
  Considere o compacto $S^{m-1}\times S^{n-1}$, como $B$ é contínua, então $B(S^{m-1}\times S^{n-1})$ é compacto.

  Definimos
  \begin{align*}
    |B|=\sup\{|B(u,v)|:u\in S^{m-1},v\in S^{n-1}\}.
  \end{align*}
  para $x\neq 0$ e $y\neq 0$, podemos escrever
  \begin{align*}
    B(x,y)&=B\left( \frac{x}{|x|}|x|,\frac{y}{|y|}|y|\right)\\
    &=B\left( \frac{x}{|x|},\frac{y}{|y|}\right)|x||y|.
  \end{align*}
  Portanto,
  \begin{align*}
    |B(x,y)|\leq |B||x||y|.
  \end{align*}
  Para $x$ ou $y$ igual a $0$ tambem se cumple.\\
  Então, sejam $x,u\in \mathbb{R}^{m}$ e $y,v\in\mathbb{R}^{n}$, logo
  \begin{align*}
    B(x+u,y+v)=B(x,y)+B(x,v)+B(u,y)+B(u,v).
  \end{align*}
  Assim
  \begin{align*}
    B(x+u,y+v)-B(x,y)=B(x,v)+B(u,y)+B(u,v).
  \end{align*}
  \textbf{Afirmação}
  \begin{align*}
    \lim_{(u,v) \to (0,0)}\frac{B(u,v)}{|(u,v)|}=0.
  \end{align*}
  Sabemos que $|B(u,v)|\leq |B||u||v|$ e $|(u,v)|\geq |u|,|v|$, então
  \begin{align*}
    \lim_{(u,v) \to (0,0)}\frac{B(u,v)}{|(u,v)|}&\leq \lim_{(u,v) \to (0,0)}\frac{|B||u||v|}{|(u,v)|}\\
    &\leq \lim_{(u,v) \to (0,0)}|B||u|=0.
  \end{align*}
  Logo $B^{\prime}(x,y)(u,v)=B(u,y)+B(x,v)$.\\
  \textbf{Afirmação} $B\in C^{\infty}$.
  \begin{itemize}
    \item Para $y$ fixo ten-se que $(x,u)\to B(u,y)$ é linear. 
    \item Para $x$ fixo ten-se que $(y,v)\to B(x,v)$ é linear. 
  \end{itemize}
  Assim, $B^{\prime}$ é uma aplicação que a cada $(x,y)$ associa uma transformação linear que depende linearmente de $x$ e $y$.\\
  Derivando novamente ao derivadas de ordem superior serão nulas.
\end{example}
\begin{example}
  Sejam $f,g:U\to\mathbb{R}^{n}$ diferenciável no ponto $a\in U\subseteq\mathbb{R}^{m}$. Então a aplicação
  \begin{align*}
    (f,g):U&\to\mathbb{R}^{n}\times\mathbb{R}^{n}\\
    x&\to(f,g)(x)=(f(x),g(x)).
  \end{align*}
  é diferenciável no ponto $a$ e a sua derivada é $(f,g)^{\prime}(a)\cdot v=(f^{\prime}(a)\cdot v,g^{\prime}(a)\cdot v)$ para tudo $v\in\mathbb{R}^{m}$.
  Além disso, $f,g\in C^{k}$, então $(f,g)\in C^{k}$.
\end{example}
